%%%%%%%%%%%%%%%%%%%%%%%%%%%%%%%%%%%%%%%%%%%%%%%%%%%%%%%%%%%%%%%%%%%%%%%%%%%%%%%%
%2345678901234567890123456789012345678901234567890123456789012345678901234567890
%        1         2         3         4         5         6         7         8

\documentclass[letterpaper, 10 pt, conference]{ieeeconf}  % Comment this line out if you need a4paper

%\documentclass[a4paper, 10pt, conference]{ieeeconf}      % Use this line for a4 paper

\IEEEoverridecommandlockouts                              % This command is only needed if 
                                                          % you want to use the \thanks command

\overrideIEEEmargins                                      % Needed to meet printer requirements.

%In case you encounter the following error:
%Error 1010 The PDF file may be corrupt (unable to open PDF file) OR
%Error 1000 An error occurred while parsing a contents stream. Unable to analyze the PDF file.
%This is a known problem with pdfLaTeX conversion filter. The file cannot be opened with acrobat reader
%Please use one of the alternatives below to circumvent this error by uncommenting one or the other
%\pdfobjcompresslevel=0
\pdfminorversion=4

% See the \addtolength command later in the file to balance the column lengths
% on the last page of the document

% The following packages can be found on http:\\www.ctan.org
\usepackage{graphics} % for pdf, bitmapped graphics files
\usepackage{epsfig} % for postscript graphics files
%\usepackage{mathptmx} % assumes new font selection scheme installed
\usepackage{times} % assumes new font selection scheme installed
\usepackage{amsmath} % assumes amsmath package installed
\usepackage{amssymb}  % assumes amsmath package installed
\usepackage[table]{xcolor}
\usepackage{mathtools,commath}
\usepackage{graphicx,import}
\usepackage{verbatim}
\usepackage{color, soul}
%\usepackage{hyperref}
\usepackage{lipsum}
\usepackage{verbatim}
\usepackage{algorithm}
\usepackage[noend]{algpseudocode}
%\usepackage{subfig}
\usepackage{multicol}
\usepackage{subfigure}
\usepackage{tikz}
\usepackage{booktabs}
\usepackage[utf8]{inputenc}
\usepackage[english]{babel}
\usepackage{todonotes}
\usepackage{diagbox}
\usepackage{color, colortbl}
\usepackage{microtype}
%\usepackage[amssymb]{SIunits}
\usepackage{siunitx} % For SI unit measures (includes SIunits package functions)
\usepackage{gensymb} % For \degree command when using also siunitx package
\usepackage{cite}
\addto\captionsenglish{\renewcommand*{\tablename}{TABLE}}

\definecolor{Gray}{gray}{0.9}
\newcolumntype{g}{>{\columncolor{Gray}}c}

\usepackage{xurl} % package to break url links
%\usepackage{breakurl}
\usepackage[breaklinks, citecolor=black]{hyperref}
\usepackage{multirow}

%% comments
\newcommand\gabrisays[1]{\textcolor[rgb]{0.00,0.45,0.74}{#1}} %blue
\newcommand\antonellosays[1]{\textcolor[rgb]{0,0.6,0}{#1}} %green
\newcommand\fabioDsays[1]{\textcolor[rgb]{0.58,0,0.82}{#1}} %purple
\newcommand\tongsays[1]{\textcolor[rgb]{0.858, 0.188, 0.478}{#1}} %pink
\newcommand\giusesays[1]{\textcolor[rgb]{0.81, 0.32, 0}{#1}} %orange
\newcommand\fabioBsays[1]{\textcolor[rgb]{0.5,0.7,0.9}{#1}} %sky-blue

%%%%%%%%%%%%%%%%%%%%%%%%%%%%%%%%%%%%%%%%%%%%%%%%%%%%%%%%%%%%%%%%%%%%%%%%%%%%%%%%

\title{\LARGE \bf
Modeling and Control of Complex-shaped VTOL Systems \\ Including Aerodynamic Forces 
}

\author{Tong Hui$^{1,2}$, Gabriele Nava$^{1}$, Giuseppe L'Erario$^{1,3}$, Fabio Di Natale$^{1}$, Antonello Paolino$^{1,4}$, \\ Fabio Bergonti$^{1,3}$, Francesco Braghin$^{2}$ and Daniele Pucci$^{1,3}$

\thanks{$^{1}$ Artificial and Mechanical Intelligence, Istituto Italiano di Tecnologia, Genova, Italy {\tt\small firstname.surname@iit.it}}%
\thanks{$^{2}$ Department of Mechanical Engineering, Politecnico di Milano, Milan, Italy {\tt\small francesco.braghin@polimi.it} {\tt\small tong.hui@mail.polimi.it}}%
\thanks{$^{3}$ School of Computer Science, Univ. of Manchester, Manchester, U.K.}
\thanks{$^{4}$ Department of Industrial Engineering, Universit$\grave{a}$ degli Studi di Napoli Federico II, Naples, Italy} 
}

\begin{document}

\maketitle
\thispagestyle{empty}
\pagestyle{empty}

%%%%%%%%%%%%%%%%%%%%%%%%%%%%%%%%%%%%%%%%%%%%%%%%%%%%%%%%%%%%%%%%%%%%%%%%%%%%%%%%
\begin{abstract}

This paper presents a modeling and control framework for complex-shaped VTOL vehicles flying in presence of non-negligible aerodynamic effects.
First, aerodynamic forces are calculated during robot flight in different operating conditions by means of Computational Fluid Dynamics (CFD) analysis.
Then, analytical models of the aerodynamics coefficients are generated from the dataset collected with CFD analysis, and used to improve the flying robot control design. We present two control strategies: compensate for the aerodynamic effects via feedback linearization and enforcing the controller robustness with gain scheduling. Simulation results on the jet-powered humanoid robot iRonCub validate the proposed approach.

\end{abstract}

%%%%%%%%%%%%%%%%%%%%%%%%%%%%%%%%%%%%%%%%%%%%%%%%%%%%%%%%%%%%%%%%%%%%%%%%%%%%%%%%

\import{text/}{introduction}
\import{text/}{background}
\import{text/}{aerodynamic_modelling}
\import{text/}{simulation_validation}
\import{text/}{conclusion}


%\addtolength{\textheight}{-12cm}   % This command serves to balance the column lengths
                                  % on the last page of the document manually. It shortens
                                  % the textheight of the last page by a suitable amount.
                                  % This command does not take effect until the next page
                                  % so it should come on the page before the last. Make
                                  % sure that you do not shorten the textheight too much.

%%%%%%%%%%%%%%%%%%%%%%%%%%%%%%%%%%%%%%%%%%%%%%%%%%%%%%%%%%%%%%%%%%%%%%%%%%%%%%%%

\import{}{biblio}

\end{document}
